\documentclass[11pt]{article}

\usepackage[a4paper,margin=2.4cm]{geometry}
\usepackage{amsmath,amssymb}
\usepackage{booktabs}
\usepackage{hyperref}
\usepackage{xcolor}

\title{Hack The World(s) -- Track 7\\Learned Cost / Value for Planning}
\author{Team HTW}
\date{\today}

\newcommand{\code}[1]{\texttt{#1}}

\begin{document}
\maketitle

\section{Goal}

The Track 7 hypothesis is that distance in latent space is a crude planning
cost. A JEPA world model can predict future latent states, but the standard MPC
objective,
\[
  J_{\mathrm{dist}}(a_{1:H}) =
  \sum_{t=1}^{H} \left\| \hat z_t - z_g \right\|_2^2,
\]
does not necessarily correlate with task success. In Two Rooms, the optimal
trajectory is often non-monotonic: the agent must first move toward the door,
which can temporarily increase Euclidean or latent distance to the final goal.

We replace this hand-crafted distance with a learned goal-conditioned value
function:
\[
  V_\psi(z_t, z_g) \in (0,1),
\]
trained on the world model's own rollouts. At planning time MPPI minimizes
\[
  J_{\mathrm{value}}(a_{1:H}) =
  \sum_{t=1}^{H} \gamma^{t-1} \left(1 - V_\psi(\hat z_t, z_g)\right).
\]
The planner therefore optimizes a quantity trained to predict reachability
rather than raw representation similarity.

\section{Why This Fits the Track}

The official track asks for three concrete changes:
\begin{itemize}
  \item a new MPC objective in \code{eb\_jepa/planning.py}, registered in
  \code{objective\_name\_map};
  \item a small value head trained alongside, or on top of, the JEPA;
  \item metrics on Two Rooms success and planning compute versus the distance
  baseline, especially on hard non-monotonic episodes.
\end{itemize}

The implementation already contains the core hook:
\code{LearnedValueMPCObjective} is registered as \code{learned\_value}, and
\code{GoalValueHead} maps \((z_t,z_g)\) to a scalar value. The next run should
make this the main experiment rather than an auxiliary maze result.

\section{Lessons from the Maze README}

The Maze README is important context, but it should not replace the Track 7
Two Rooms experiment. It teaches four design constraints:

\begin{itemize}
  \item \textbf{Freeze a proven world model first.} Maze co-training lowered
  subgoal MSE but hurt success, because moving the encoder degraded the
  wall-aware dynamics. Our main Track 7 comparison therefore trains the value
  head on a frozen WM before attempting any joint fine-tuning.
  \item \textbf{Keep the comparison controlled.} The Maze value result is strong
  only when the planner scaffolding is matched: same WM, same planner, same
  episodes, only the cost changes. We copy exactly that protocol for Two Rooms.
  \item \textbf{Do not hide oracle guidance in the planner.} Maze A* waypoints
  are useful as a stress-test, but the Track 7 story should not depend on A*.
  A* can be an oracle for analysis, not the source of the main action signal.
  \item \textbf{Horizon matters.} Maze greedy-global planning stays at 0\%
  because a value trained on short windows cannot guide long-range navigation.
  Two Rooms is the right first target because its non-monotonic door crossing is
  hard enough to expose the flaw of latent distance while remaining within the
  value's rollout horizon.
\end{itemize}

Thus the main Track 7 run is not ``Maze hierarchy again''. It is the cleaner
learned-cost experiment requested by the track, with Maze kept as supporting
evidence and an ambitious extension.

\section{Training Objective}

We train a value head with a TD-MPC-style bootstrap. For each sampled trajectory
window, the hindsight goal is the final encoded frame \(z_g=z_T\). The target is
\[
  y_t = r_t + \gamma (1-d_t) V_{\bar\psi}(z_{t+1}, z_g),
\]
where \(V_{\bar\psi}\) is an EMA target network, \(r_t=1\) only at the endpoint,
and \(d_t\) marks the terminal endpoint transition.

The loss is computed on both encoded real latents and imagined world-model
rollout latents:
\[
  \mathcal L_{\mathrm{value}} =
  \left\|V_\psi(z_t,z_g)-y_t\right\|_2^2
  +
  \left\|V_\psi(\hat z_t,z_g)-y_t\right\|_2^2.
\]
The second term matters: MPPI will query the value on imagined states, so the
head must be calibrated on the distribution produced by the world model.

\section{Controlled Comparison}

The main experiment changes exactly one thing: the planning objective.
\[
\text{same world model, same MPPI, same seeds, same episodes}
\]
\[
\code{repr\_dist} \quad \mathrm{vs.} \quad \code{learned\_value}.
\]

\begin{table}[h]
\centering
\begin{tabular}{lll}
\toprule
Condition & Objective & Purpose \\
\midrule
Baseline & \code{repr\_dist} & Official distance-in-latent-space baseline \\
Ours & \code{learned\_value} & TD-MPC value learned from WM rollouts \\
Oracle report only & A* / shortest path & Upper bound, not used for training in main result \\
\bottomrule
\end{tabular}
\caption{Track 7 comparison. A* is only an analysis oracle in the main story.}
\end{table}

\section{Metrics}

We report:
\begin{itemize}
  \item success rate on Two Rooms;
  \item mean final distance to goal;
  \item average episode wall time and approximate world-model calls;
  \item success on the hard subset: start and goal on opposite sides of the wall,
  normal-level layout, long door-seeking trajectories;
  \item qualitative GIFs: one success, one failure, one non-monotonic case.
\end{itemize}

For a compute story, we run at least two MPPI budgets:
\[
  \code{num\_samples} \in \{50, 200\}, \qquad
  \code{n\_iters}=20, \qquad
  \code{plan\_length}=90.
\]
This gives a small Pareto curve: success versus planning time.

\section{Expected Jury Story}

The concise claim should be:

\begin{quote}
The baseline world-model planner optimizes latent similarity, but latent
similarity is not a task cost. We train a TD-MPC-style value head on the JEPA
world model's own rollouts and use it as the MPPI objective. On non-monotonic
Two Rooms episodes, this value provides a reachability signal that is better
aligned with success than raw latent distance.
\end{quote}

\section{Execution Checklist}

\begin{enumerate}
  \item Train or reuse a stable Two Rooms world model checkpoint.
  \item Train \code{GoalValueHead} on the frozen world model using
  \code{train\_value.yaml}.
  \item Evaluate \code{repr\_dist} and \code{learned\_value} on the same
  \code{eval\_track7.yaml} episodes.
  \item Repeat at low and standard MPPI compute budgets.
  \item Export table, Pareto plot, and 2--3 GIFs.
  \item If time remains, run an ambitious joint fine-tune with
  \code{freeze\_world\_model=false} and small \code{value\_coeff}.
\end{enumerate}

\section{Risk Management}

If the learned value underperforms globally, we still have a useful result:
the value can be analyzed by horizon. A short-window TD target may improve
near-goal decisions while failing on long-range navigation. This is a clean
negative or partial result and directly supports the paper's stated limitation:
latent-distance costs are not enough, but learned values need horizon-aware
training.

\end{document}
